Deep learning has transformed pattern recognition by enabling end-to-end learning of hierarchical representations~\cite{goodfellow2016deep}. In physiological signal analysis, these paradigms enable decoding of complex biological time-series like EEG, traditionally constrained by handcrafted features.

Brain--Computer Interfaces (BCIs) provide direct communication between neural activity and external devices, offering potential to restore independence for individuals with severe motor impairments such as ALS or spinal cord injuries~\cite{lu2017deep, luo2018lstm}. Motor Imagery (MI)—mental simulation of movement—activates sensorimotor cortical regions comparably to actual movement, making it suitable for intuitive control of prosthetics and wheelchairs~\cite{xu2018wavelet}.

However, MI-EEG decoding remains challenging due to non-stationarity, low signal-to-noise ratios, and inter-subject variability~\cite{autthasan2021min2net}. Conventional approaches like FBCSP require domain expertise and generalize poorly~\cite{xu2019deep}. While deep learning addresses some limitations, limited evidence exists on whether complex architectures consistently outperform lightweight models in low-data BCI scenarios.

This study comparatively analyzes two end-to-end architectures for three-class MI-EEG classification (left-hand, right-hand, both-feet): EEGNet, a parameter-efficient CNN, and ShallowConvNet, emphasizing interpretable spectral--spatial filtering. Evaluating both under identical conditions, this work investigates whether increased complexity yields measurable gains, with attention to limited-data settings and cross-subject generalization. Beyond accuracy, error analysis identifies failure factors from modeling (underfitting, overfitting) and data-centric (noise, variability) perspectives.

This study contributes: (1) an end-to-end EEG classification pipeline with preprocessing, filtering, artifact mitigation, and rigorous validation; and (2) systematic comparison between EEGNet and ShallowConvNet focusing on performance, generalization, and error characteristics.
